\documentclass[conference]{IEEEtran}
\IEEEoverridecommandlockouts
\usepackage{cite}
\usepackage{amsmath,amssymb,amsfonts}
\usepackage{algorithmic}
\usepackage{algorithm}
\usepackage{graphicx}
\usepackage{textcomp}
\usepackage{xcolor}
\usepackage{booktabs}
\usepackage{multirow}
\usepackage{subcaption}
\usepackage{hyperref}

\def\BibTeX{{\rm B\kern-.05em{\sc i\kern-.025em b}\kern-.08em
    T\kern-.1667em\lower.7ex\hbox{E}\kern-.125emX}}

\begin{document}
\title{Not All Items Are Created Equal\\
\vspace{0.5cm}
\large Importance-Aware Next Basket Recommendation}

\author{
\IEEEauthorblockN{Dhruv Ramesh Joshi}
\IEEEauthorblockA{\textit{IIIT-Bangalore}}
\and
\IEEEauthorblockN{R Ricky Roger}
\IEEEauthorblockA{\textit{IIIT-Bangalore}}
\and
\IEEEauthorblockN{Siddharth Prashant Kini}
\IEEEauthorblockA{\textit{IIIT-Bangalore}}
\and
\IEEEauthorblockN{Subhash Hari}
\IEEEauthorblockA{\textit{IIIT-Bangalore}}
}

\maketitle

%-----------------------------------------------------------------------
\begin{abstract}
Next basket recommendation predicts the full set of items a user will
purchase next given their history. Existing methods encode each basket
as a flat vector and treat every item as an equally important signal of
user intent. We argue this is structurally wrong: every basket is
organized around a small number of intent-defining items, with the
remainder serving as peripheral fill. We propose an architecture that
makes this distinction explicit at every stage. A BERT-style intra-basket
encoder produces per-item contextualized representations, from which a
learned importance head derives per-item importance weights initialized
from a geometric leave-one-out attribution score corrected for
population frequency. An elementwise-gated dual-stream representation
blends a full-basket summary with an importance-weighted intent centroid.
A causal GPT with rotary position embeddings models the temporal sequence
of basket representations and produces a single predicted next-basket
vector. Decoding proceeds in two conditioned stages: core items are
predicted first via a low-rank orthogonal intent projection, and
peripheral items are then predicted from a fill head whose query is
explicitly conditioned on the discovered core, bridged by a
differentiable soft intent context during training.
We evaluate on TaFeng and Instacart benchmarks and demonstrate
improvements over frequency baselines and ablated variants,
validating the importance-aware design.
\end{abstract}

\begin{IEEEkeywords}
next basket recommendation, purchase intent, importance scoring,
transformer, hierarchical decoding, sequential recommendation
\end{IEEEkeywords}


%-----------------------------------------------------------------------
\section{Introduction}
Next Basket Recommendation (NBR) is a well-established problem in the recommender systems literature. Given a sequence of historical baskets---sets of items purchased together---the goal is to predict the set of items a user will select in their next interaction. This setting is prominent in e-commerce and grocery shopping, where users routinely purchase multiple items in a single session.

A user who buys chicken, basmati rice, saffron, yogurt, and cooking oil
in one trip is not selecting five independent items. They are executing
a meal plan. The basket is the physical realization of a latent intent.
Remove the chicken and the basket loses its organizing purpose. Remove
the cooking oil and the basket remains a coherent meal.

This asymmetry is invisible to every existing NBR system. Standard methods encode each basket as a mean or max pool over item embeddings and treat all items as equal contributors. The model receives identical gradient signal whether it mispredicts the meal anchor or an incidental staple.

We propose an architecture grounded in the observation that a basket is
not a flat bag but a hierarchical satisfaction of a purchase intent.
Concretely, we contribute:
\begin{itemize}
\item A \textbf{geometric importance score} per item, measuring how much
  the basket representation shifts when that item is removed, corrected
  for population frequency.
\item A \textbf{parallel importance head} that produces per-item importance weights, initialized from the geometric score and refined during training.
\item A \textbf{dual-stream gated basket representation} that blends the
  full-basket CLS summary with an importance-weighted intent centroid.
\item A \textbf{conditioned two-stage decoder}: core items are predicted
  from a low-rank intent projection; fill items are
  then predicted from a fill head conditioned on the discovered core via a differentiable soft bridge.
\end{itemize}


%-----------------------------------------------------------------------
\section{Related Work}
\label{sec:related}

\subsection{Frequency and Neighbor Methods}
The simplest NBR baselines rely on item frequency.
P-TopFreq and G-TopFreq recommend the most personally or globally
frequent items respectively. Li et al.~\cite{li2023next} propose
GP-TopFreq as a strong hybrid baseline. TIFUKNN~\cite{hu2020modeling} and
UP-CF@r~\cite{faggioli2020recency} add temporal frequency dynamics and
recency weighting to collaborative filtering.

\subsection{RNN and Graph Methods}
DREAM~\cite{yu2016dynamic} introduced RNN-based temporal modeling of
basket sequences. HRM~\cite{wang2015learning} formalized the two-level
item-within-basket and basket-within-sequence hierarchy.
DNNTSP~\cite{yu2020predicting} built a dynamic co-occurrence graph.
Sets2Sets~\cite{hu2019sets2sets} encodes frequency as an explicit input feature.
Beacon~\cite{le2019correlation} captures pairwise item correlations.

\subsection{Contrastive Methods}
CLEA~\cite{qin2021world} learns a binary relevance mask via Gumbel-Softmax to separate relevant from irrelevant basket items. Our importance score is a continuous geometric measurement computed independently of any task objective.
CL4NBR~\cite{zhu2025contrastive} adds contrastive augmentation over basket sequences with uniform item treatment.

\subsection{Transformer Methods}
SASRec~\cite{kang2018self} demonstrated that causal self-attention
outperforms RNNs for sequential recommendation.
BERT4Rec~\cite{sun2019bert4rec} applied bidirectional masked attention.
Our Layer 1 and Layer 2 encoders are directly inspired by these two papers. NEST is the closest architectural ancestor: a set-within-encounter bidirectional encoder followed by a causal sequence encoder with RoPE. We extend this with importance scoring, dual-stream fusion, and conditioned two-stage decoding.

\subsection{The Repeat-Explore Problem}
ReCANet~\cite{ariannezhad2022recanet} established that over 54\% of NBR recall comes from repeat items. Li et al.~\cite{li2023next} showed no published NBR method properly balances repeat and explore prediction, introducing separate Repeat Recall@$K$ and Explore Recall@$K$ metrics. Our IDF-corrected importance score directly addresses this.


%-----------------------------------------------------------------------
\section{Methodology}
\label{sec:method}

\subsection{Item Embedding Initialization}
We initialize the shared item embedding matrix
$E \in \mathbb{R}^{|V| \times d}$ via skip-gram Word2Vec on the basket
corpus, treating each basket as a sentence and each item as a token.
$E$ is fully trainable throughout all phases.

\subsection{Intra-Basket Encoder with Importance Head}
Each basket $B_t = \{i_1, \ldots, i_K\}$ is encoded by a BERT-style
transformer. A learned CLS token is prepended:
\begin{equation}
X = [e_{\text{CLS}},\ E[i_1],\ \ldots,\ E[i_K]] \in \mathbb{R}^{(K+1) \times d}
\end{equation}

We apply $L_1$ layers of full bidirectional self-attention with no
positional encoding (baskets are unordered sets) and no causal mask.
The output is $[h_{\text{CLS}}, h_1, \ldots, h_K]$.

A parallel importance head maps each contextualized representation to an importance weight:
\begin{equation}
\alpha_i = \sigma\!\left(W_2\, \operatorname{GELU}(W_1 h_i + b_1) + b_2\right) \in [0,1]
\end{equation}

An auxiliary masked language modeling loss provides a direct training signal for the encoder:
\begin{equation}
\mathcal{L}_{\text{MLM}} = \operatorname{CE}\!\left(E\, h_k^{\text{masked}},\ i_k^{\text{true}}\right)
\end{equation}

\subsection{Importance Score Initialization}
After a warmup phase, we freeze the BERT encoder and compute a per-item importance score via leave-one-out CLS shift:
\begin{equation}
\delta(i, B_t) = \|h_{\text{CLS}}(B_t) - h_{\text{CLS}}(B_t \setminus \{i\})\|_2
\end{equation}

We normalize within each basket and aggregate across appearances:
\begin{align}
\hat{\delta}(i, B_t) &= \frac{\delta(i, B_t)}{\sum_{j \in B_t} \delta(j, B_t)} \\
\bar{\delta}_i        &= \frac{1}{|\{(u,t):i \in B_t\}|} \sum_{(u,t):\,i \in B_t} \hat{\delta}(i, B_t)
\end{align}

An IDF correction separates structural importance from frequency:
\begin{equation}
\alpha_i^{IDF} = \bar{\delta}_i \cdot \log\frac{N}{df_i}
\end{equation}

The importance head is pre-trained to reproduce these scores via MSE before main training.

\subsection{Dual-Stream Gated Basket Representation}
Two basket-level representations are constructed:
\begin{align}
b_t^{\text{full}} &= h_{\text{CLS}} \\
b_t^{\text{core}} &= \frac{\sum_k \alpha_{i_k}\, h_k}{\sum_k \alpha_{i_k}}
\end{align}

An elementwise gate fuses them:
\begin{align}
g_t &= \sigma\!\left(W_g [b_t^{\text{full}};\, b_t^{\text{core}}]\right) \in [0,1]^d \\
b_t &= g_t \odot b_t^{\text{full}} + (1 - g_t) \odot b_t^{\text{core}}
\end{align}

\subsection{Inter-Basket Causal GPT}
The sequence $(b_1, \ldots, b_T)$ is fed into a causal GPT with $L_2$ transformer layers and rotary position embeddings (RoPE) applied to Q and K:
\begin{equation}
A = \operatorname{softmax}\!\left(\frac{\operatorname{RoPE}(Q)\,\operatorname{RoPE}(K)^\top}{\sqrt{d}} + M_{\text{causal}}\right)
\end{equation}

The output at the final position is $\hat{h}_{T+1}$, the predicted next-basket representation.

\subsection{Orthogonal Intent-Fill Decomposition}
A learned low-rank orthogonal projection $P \in \mathbb{R}^{d \times d_k}$ ($P^\top P = I_{d_k}$) separates intent from fill:
\begin{align}
\hat{h}^{\text{intent}} &= PP^\top \hat{h}_{T+1} \\
\hat{h}^{\text{fill}}   &= \hat{h}_{T+1} - PP^\top \hat{h}_{T+1}
\end{align}

Orthonormality is enforced via periodic Gram-Schmidt re-orthonormalization and a regularizer $\mathcal{L}_{\text{orth}} = \|\hat{h}^{\text{intent}} \cdot \hat{h}^{\text{fill}}\|^2$.

\subsection{Conditioned Two-Stage Decoding}

\paragraph{Stage 1 --- Core item prediction.}
Items scored against the intent component: $s_i^{\text{intent}} = e_i \cdot \hat{h}^{\text{intent}}$.
A soft intent context is computed during training:
\begin{equation}
\tilde{c} = \sum_{i \in V} \operatorname{softmax}_\tau(s^{\text{intent}})_i \cdot e_i
\end{equation}

At inference, a residual loop selects $K_1$ core items:
\begin{align}
\hat{i}_k &= \operatorname*{arg\,max}_{i \notin \hat{C}}\; e_i \cdot r_{k-1}^{\text{intent}} \\
r_k^{\text{intent}} &= r_{k-1}^{\text{intent}} - PP^\top e_{\hat{i}_k}
\end{align}

\paragraph{Stage 2 --- Conditioned fill prediction.}
The fill query is shifted by the intent context:
\begin{equation}
\hat{h}^{\text{fill} \mid \text{intent}} = \operatorname{LN}\!\left(\hat{h}^{\text{fill}} + W_{\text{cond}}\, \tilde{c}\right)
\end{equation}

Fill items are decoded via a residual loop in the fill subspace:
\begin{align}
\hat{i}_{K_1+k} &= \operatorname*{arg\,max}_{i \notin \hat{C}}\; e_i \cdot r_{k-1}^{\text{fill}} \\
r_k^{\text{fill}} &= r_{k-1}^{\text{fill}} - (I - PP^\top) e_{\hat{i}_{K_1+k}}
\end{align}

\subsection{Training Objective}
Targets are partitioned using threshold $\tau_\alpha$ on $\alpha^{IDF}$:
\begin{equation}
C^*_{T+1} = \{i \in B_{T+1} : \alpha_i^{IDF} > \tau_\alpha\},\quad
F^*_{T+1} = B_{T+1} \setminus C^*_{T+1}
\end{equation}

The total loss:
\begin{equation}
\mathcal{L} = \mathcal{L}_{\text{intent}}
            + \lambda\,\mathcal{L}_{\text{fill}}
            + \gamma\,\mathcal{L}_{\text{orth}}
            + \eta\,\mathcal{L}_{\text{MLM}}
\end{equation}

where $\mathcal{L}_{\text{intent}}$ is $\alpha^{IDF}$-weighted BCE on core items, $\mathcal{L}_{\text{fill}}$ is uniform BCE on fill items.

\subsection{Training Schedule}
Training proceeds in four phases:
\textbf{Phase 0}: $E$ initialized via skip-gram.
\textbf{Phase 1}: BERT encoder trained with MLM.
\textbf{Phase 2}: BERT frozen; $\alpha^{IDF}$ computed via leave-one-out; importance head pre-trained.
\textbf{Phase 3}: All components unfrozen; full loss with validation every 3 epochs and best-model checkpointing.


%-----------------------------------------------------------------------
\section{Experimental Setup}
\label{sec:experiments}

\subsection{Datasets}

We evaluate on two standard grocery benchmarks commonly used in NBR literature.

\begin{table}[h]
\centering
\caption{Dataset statistics after preprocessing (min 3 baskets/user, min 5 item frequency).}
\label{tab:datasets}
\begin{tabular}{lrrrr}
\toprule
\textbf{Dataset} & \textbf{Users} & \textbf{Items} & \textbf{Baskets} & \textbf{Avg. Size} \\
\midrule
%% TODO: Fill in actual dataset statistics after preprocessing
TaFeng     & --- & --- & --- & --- \\
Instacart  & --- & --- & --- & --- \\
\bottomrule
\end{tabular}
\end{table}

\subsection{Baselines}
We compare against the following baselines:
\begin{itemize}
\item \textbf{G-TopFreq}: Recommends globally most frequent items.
\item \textbf{P-TopFreq}: Recommends personally most frequent items per user.
\item \textbf{GP-TopFreq}: Hybrid of global and personal frequency~\cite{li2023next}.
\item \textbf{Dual-Stream (no two-stage)}: Our encoder-GPT pipeline with standard BCE loss and no intent-fill decomposition.
\end{itemize}

\subsection{Evaluation Metrics}
Following Li et al.~\cite{li2023next}, we report:
\begin{itemize}
\item \textbf{Recall@$K$}: Fraction of ground-truth items appearing in top-$K$ predictions.
\item \textbf{NDCG@$K$}: Normalized discounted cumulative gain at $K$.
\end{itemize}
We evaluate at $K \in \{10, 20\}$.

\subsection{Implementation Details}
All models use embedding dimension $d = 128$, intent dimension $d_k = 32$, $L_1 = 2$ BERT layers, $L_2 = 2$ GPT layers, 4 attention heads, dropout $0.22$, and are trained with AdamW ($\text{lr}=10^{-3}$, weight decay $4 \times 10^{-3}$). Phase 3 uses $K_1 = 2$ core items and $K_2 = 18$ fill items. Loss weights are $\lambda = 1.0$, $\gamma = 0.1$, $\eta = 0.3$. Gram-Schmidt re-orthonormalization is applied every 100 steps. Training is run for 20 epochs with validation every 3 epochs. All experiments are run on a single NVIDIA GPU with PyTorch.


%-----------------------------------------------------------------------
\section{Results}
\label{sec:results}

\subsection{Main Results}

\begin{table*}[t]
\centering
\caption{Next basket recommendation performance on TaFeng and Instacart. Best results are \textbf{bolded}. All metrics are averaged over the test set.}
\label{tab:main_results}
\begin{tabular}{l|cccc|cccc}
\toprule
& \multicolumn{4}{c|}{\textbf{TaFeng}} & \multicolumn{4}{c}{\textbf{Instacart}} \\
\textbf{Method} & R@10 & R@20 & N@10 & N@20 & R@10 & R@20 & N@10 & N@20 \\
\midrule
G-TopFreq       & --- & --- & --- & --- & --- & --- & --- & --- \\
P-TopFreq       & --- & --- & --- & --- & --- & --- & --- & --- \\
GP-TopFreq      & --- & --- & --- & --- & --- & --- & --- & --- \\
\midrule
Dual-Stream     & --- & --- & --- & --- & --- & --- & --- & --- \\
\textbf{Full Model (Ours)} & --- & --- & --- & --- & --- & --- & --- & --- \\
\bottomrule
\end{tabular}
\end{table*}

%% TODO: Fill Table~\ref{tab:main_results} with results from:
%%   - scripts/evaluate_baselines.py (frequency baselines)
%%   - experiments/dual_stream.py (dual-stream ablation)
%%   - experiments/full_model.py (full model)
%% Replace --- with actual numbers.

Table~\ref{tab:main_results} presents the main results.
%% TODO: Add 2-3 sentences of analysis once numbers are available.
%% Expected narrative: Full model should outperform baselines, especially
%% on explore items. Dual-stream (without two-stage decode) should be
%% competitive but weaker on fill item prediction.


\subsection{Training Dynamics}

%% TODO: Add W&B training curves as figures. Export from W&B dashboard.
%% Suggested plots:
%%   - Figure: Training loss over steps (total, intent, fill, orth, mlm)
%%   - Figure: Validation Recall@10 over epochs

\begin{figure}[h]
\centering
%% \includegraphics[width=\columnwidth]{figures/training_loss.png}
\fbox{\parbox{0.9\columnwidth}{\centering\vspace{2cm}\textit{TODO: Training loss curves from W\&B}\vspace{2cm}}}
\caption{Phase 3 training loss decomposition over steps. The intent and fill losses decrease steadily while the orthogonality loss remains near zero, confirming that $P$ maintains near-orthonormal columns.}
\label{fig:training_loss}
\end{figure}

\begin{figure}[h]
\centering
%% \includegraphics[width=\columnwidth]{figures/val_recall.png}
\fbox{\parbox{0.9\columnwidth}{\centering\vspace{2cm}\textit{TODO: Validation Recall@10 over epochs from W\&B}\vspace{2cm}}}
\caption{Validation Recall@10 evaluated every 3 epochs during Phase 3 training.}
\label{fig:val_recall}
\end{figure}


%-----------------------------------------------------------------------
\section{Analysis}
\label{sec:analysis}

\subsection{Ablation Study}

To validate each architectural component, we ablate key design decisions.

\begin{table}[h]
\centering
\caption{Ablation study on TaFeng. Each row removes one component from the full model.}
\label{tab:ablation}
\begin{tabular}{lcccc}
\toprule
\textbf{Variant} & R@10 & R@20 & N@10 & N@20 \\
\midrule
Full Model                    & --- & --- & --- & --- \\
\midrule
$-$ Importance Head           & --- & --- & --- & --- \\
$-$ Dual-Stream Gate          & --- & --- & --- & --- \\
$-$ Two-Stage Decode          & --- & --- & --- & --- \\
$-$ IDF Correction            & --- & --- & --- & --- \\
$-$ MLM Auxiliary Loss        & --- & --- & --- & --- \\
$-$ Gram-Schmidt Re-orth      & --- & --- & --- & --- \\
\bottomrule
\end{tabular}
\end{table}

%% TODO: Run ablation experiments and fill Table~\ref{tab:ablation}.
%% For each variant, train the full model with one component disabled:
%%   - No importance head: set all alpha_i = 1.0 (uniform)
%%   - No dual-stream gate: use only CLS repr (skip fusion)
%%   - No two-stage decode: single BCE head, no projection
%%   - No IDF: use raw geometric delta without log(N/df) correction
%%   - No MLM: set eta=0 in loss weights
%%   - No re-orth: disable orthogonalize_projection() calls

%% TODO: Add 2-3 sentences per ablation once numbers are available.

\subsection{Dual-Stream Gate Analysis}

\begin{table}[h]
\centering
\caption{Average gate value $\bar{g}$ across validation baskets, measuring the model's learned reliance on full-basket CLS ($g \to 1$) versus importance-weighted core ($g \to 0$).}
\label{tab:gate}
\begin{tabular}{lc}
\toprule
\textbf{Dataset} & $\bar{g}$ \\
\midrule
%% TODO: Extract from dual_stream.py gate analysis output
TaFeng    & --- \\
Instacart & --- \\
\bottomrule
\end{tabular}
\end{table}

%% TODO: Interpret the gate value. If g > 0.5, the model relies more on
%% full-basket CLS; if g < 0.5, it leans toward importance-weighted core.
%% Discuss what this means for the importance head's contribution.

\subsection{Importance Score Distribution}

%% TODO: Add a histogram or violin plot of alpha_IDF scores.
%% This can be generated from importance_scores.npz.

\begin{figure}[h]
\centering
%% \includegraphics[width=\columnwidth]{figures/importance_dist.png}
\fbox{\parbox{0.9\columnwidth}{\centering\vspace{2cm}\textit{TODO: Distribution of $\alpha^{IDF}$ scores}\vspace{2cm}}}
\caption{Distribution of geometric importance scores $\alpha^{IDF}$ across all items. The bimodal distribution supports our hypothesis that items cluster into load-bearing (high $\alpha$) and peripheral (low $\alpha$) categories.}
\label{fig:importance_dist}
\end{figure}

\subsection{Effect of $K_1$ and $K_2$}

\begin{table}[h]
\centering
\caption{Sensitivity to core/fill split $K_1 + K_2 = 20$.}
\label{tab:k1k2}
\begin{tabular}{cc|cc}
\toprule
$K_1$ & $K_2$ & R@20 & N@20 \\
\midrule
%% TODO: Run inference with different K1/K2 splits
1  & 19 & --- & --- \\
2  & 18 & --- & --- \\
3  & 17 & --- & --- \\
5  & 15 & --- & --- \\
10 & 10 & --- & --- \\
\bottomrule
\end{tabular}
\end{table}

%% TODO: Discuss optimal K1/K2 split. Expected: K1=2-3 is optimal
%% because most baskets have only 1-3 truly intent-defining items.

\subsection{Qualitative Examples}

%% TODO: Pick 2-3 example users from the test set and show:
%%   - Their purchase history (last 3 baskets)
%%   - The model's predicted core items (Stage 1)
%%   - The model's predicted fill items (Stage 2)
%%   - Ground truth next basket
%% This illustrates the two-stage decoding in action.


%-----------------------------------------------------------------------
\section{Discussion}
\label{sec:discussion}

%% TODO: Expand after results are available. Key points to cover:

\paragraph{Strengths.}
The importance-aware architecture offers three structural advantages over flat NBR methods. First, the geometric importance score provides a principled, task-independent measurement of item contribution that does not conflate frequency with structural importance. Second, the dual-stream gate learns to adaptively balance full-basket and importance-filtered representations as training progresses. Third, the two-stage conditioned decoder prevents intent items from crowding fill slots and ensures fill predictions are coherent with the discovered intent.

\paragraph{Limitations.}
The current system has several limitations.
The leave-one-out importance computation (Phase 2) requires $O(N \times \bar{S})$ forward passes where $\bar{S}$ is average basket size, which is expensive for large datasets.
The fixed threshold $\tau_\alpha$ for core/fill partitioning may not generalize across datasets.
The soft-to-hard transition between training ($\tilde{c}$ via softmax) and inference (mean of $K_1$ embeddings) introduces a train-test mismatch that could be addressed via temperature annealing or Gumbel-Softmax.

\paragraph{Computational Cost.}
%% TODO: Add actual training times and inference latency
Phase 1 (BERT warmup): --- hours.
Phase 2 (importance computation): --- hours.
Phase 3 (joint training): --- hours per epoch.
Inference: --- ms per user.


%-----------------------------------------------------------------------
\section{Conclusion}
\label{sec:conclusion}

We presented an intent-aware next basket recommendation system that
makes explicit the asymmetry between load-bearing and peripheral items.
A geometric importance score, corrected for population frequency, initializes a learned importance head that shapes both the basket representation and the training signal. A dual-stream gated encoder blends full-basket context with an importance-filtered intent centroid. A causal GPT produces a predicted next-basket vector, which a low-rank orthogonal projection decomposes into intent and fill components. Decoding proceeds in two conditioned stages connected by a differentiable soft intent context.

%% TODO: Add 1-2 sentences summarizing key empirical findings once results are available.
%% e.g., "On TaFeng, our full model achieves X\% Recall@10, outperforming
%% GP-TopFreq by Y\% and the dual-stream ablation by Z\%."

Future work will extend evaluation to Dunnhumby with separate Repeat Recall@$K$ and Explore Recall@$K$ metrics, explore temperature annealing to bridge the soft-hard gap, and investigate finer hierarchical intent structures.

%-----------------------------------------------------------------------
\bibliographystyle{IEEEtran}
\bibliography{refs}

\end{document}
